
\section{Bối cảnh và Động lực nghiên cứu}

Sự phát triển mạnh mẽ của các mô hình ngôn ngữ lớn (LLMs – Large Language Models) như GPT, Claude và Gemini trong những năm gần đây đã tạo ra những thay đổi đáng kể trong lĩnh vực Trí tuệ Nhân tạo. Các mô hình này thể hiện khả năng vượt trội trong việc hiểu và sinh ngôn ngữ tự nhiên, từ đó được ứng dụng rộng rãi trong nhiều bài toán như hỗ trợ khách hàng, tổng hợp thông tin hay tạo nội dung tự động. Tuy nhiên, việc đưa LLMs vào vận hành trong các hệ thống thực tế vẫn còn gặp nhiều hạn chế, đặc biệt là khả năng truy cập và khai thác dữ liệu chuyên ngành một cách chính xác và có kiểm chứng.

Trong bối cảnh đó, bài toán **trả lời câu hỏi lịch sử** đặt ra yêu cầu đặc thù: mô hình không chỉ cần hiểu ngôn ngữ mà còn phải truy xuất được các sự kiện, mốc thời gian và bối cảnh lịch sử chính xác. Do dữ liệu lịch sử thường có cấu trúc phức tạp, phân tán và đòi hỏi tính chính xác cao, các mô hình LLM truyền thống dễ gặp phải hiện tượng “ảo giác” (hallucination), đưa ra thông tin sai lệch hoặc suy luận thiếu cơ sở khi chỉ dựa vào dữ liệu huấn luyện.

Để khắc phục những hạn chế này, các phương pháp kết hợp giữa LLM và dữ liệu bên ngoài đã được phát triển, trong đó đáng chú ý là **Retrieval-Augmented Generation (RAG)**. Tuy nhiên, với dữ liệu lịch sử có nhiều mối quan hệ phức tạp giữa các nhân vật, sự kiện và thời kỳ, RAG truyền thống – vốn dựa trên truy vấn tài liệu dạng văn bản – gặp hạn chế trong việc mô hình hóa và truy xuất tri thức có cấu trúc.

Điều này dẫn đến sự xuất hiện của **GraphRAG**, một phương pháp kết hợp khả năng sinh ngôn ngữ của mô hình LLM với sức mạnh biểu diễn tri thức của đồ thị. GraphRAG cho phép mô hình truy cập vào các mạng lưới quan hệ lịch sử được biểu diễn dưới dạng đồ thị tri thức (knowledge graph), từ đó hỗ trợ quá trình suy luận, đối chiếu và trả lời câu hỏi chính xác hơn. Bằng cách sử dụng đồ thị để lưu trữ các quan hệ như *nhân vật – sự kiện*, *nguyên nhân – hệ quả*, *địa điểm – thời gian*, GraphRAG giúp LLM tránh được sai lệch ngữ nghĩa và đảm bảo tính mạch lạc trong câu trả lời.

Sự kết hợp giữa LLM và GraphRAG cũng tận dụng được khả năng ghi nhớ ngữ cảnh, suy luận theo chuỗi và lập kế hoạch của các mô hình hiện đại. Điều này giúp mô hình không chỉ trả lời câu hỏi đơn lẻ mà còn xử lý các truy vấn phức tạp, đòi hỏi tổng hợp và phân tích nhiều lớp thông tin lịch sử. Nhờ đó, các hệ thống AI có thể đưa ra câu trả lời chính xác hơn, nhất quán hơn và phù hợp với các nguồn dữ liệu đã được kiểm chứng.

Trong bối cảnh nhu cầu học tập, nghiên cứu lịch sử và xây dựng hệ thống hỗ trợ giáo dục ngày càng tăng, việc ứng dụng LLM kết hợp GraphRAG vào bài toán trả lời câu hỏi lịch sử trở thành một hướng tiếp cận đầy tiềm năng. Đây cũng chính là động lực cho nghiên cứu trong đề tài này, nhằm xây dựng một hệ thống có khả năng khai thác tri thức lịch sử một cách hiệu quả, đáng tin cậy và phù hợp với yêu cầu thực tiễn.

\section{Tổng quan vấn đề}
Việc triển khai các hệ thống AI sử dụng mô hình ngôn ngữ lớn kết hợp phương pháp GraphRAG đặt ra nhiều thách thức, ảnh hưởng đáng kể đến khả năng mở rộng, độ tin cậy và khả năng ứng dụng trong thực tế. Mặc dù GraphRAG mang lại cách tiếp cận linh hoạt và có cấu trúc để liên kết mô hình LLM với các nguồn đồ thị tri thức, quá trình triển khai một hệ thống hoàn chỉnh vẫn còn phức tạp, đòi hỏi nhiều thao tác thủ công và phụ thuộc vào kinh nghiệm chuyên sâu của người phát triển.

Một trong những khó khăn lớn xuất phát từ việc lựa chọn, xây dựng và tích hợp các thành phần của hệ thống đồ thị tri thức. Với bài toán lịch sử, các nguồn dữ liệu thường nằm rải rác trong nhiều kho tư liệu, sách, văn bản, hoặc cơ sở dữ liệu mở. Việc chuẩn hóa dữ liệu, xây dựng đồ thị tri thức lịch sử và kết nối chúng với mô hình LLM thông qua GraphRAG đòi hỏi quy trình nhiều bước, bao gồm thu thập dữ liệu, trích xuất thực thể, phân tích quan hệ, và tối ưu hóa cấu trúc đồ thị. Nếu không có quy trình triển khai thống nhất và tự động hóa, người phát triển dễ gặp sai sót trong cấu hình, khó đảm bảo nhất quán giữa các phiên bản dữ liệu và phải tốn nhiều thời gian để tinh chỉnh hệ thống.

Bên cạnh đó, **khả năng mở rộng** (scalability) cũng là một vấn đề quan trọng khi các hệ thống trả lời câu hỏi ngày càng được sử dụng cho nhiều lượng truy vấn hơn và yêu cầu phản hồi nhanh hơn. Trong các tình huống thực tế—chẳng hạn như hỗ trợ học tập lịch sử, tra cứu sự kiện theo yêu cầu, hoặc khai thác tri thức trong môi trường giáo dục trực tuyến—hệ thống phải xử lý nhiều loại câu hỏi khác nhau và truy xuất cùng lúc nhiều nguồn dữ liệu từ đồ thị tri thức. Các mô hình GraphRAG nếu không được thiết kế tối ưu về hạ tầng có thể gặp khó khăn trong việc đáp ứng độ trễ thấp, xử lý song song hiệu quả hoặc phân phối tài nguyên theo nhu cầu tăng dần của người dùng.

Một thách thức khác là **giám sát và quan sát hệ thống** (monitoring \& observability). Do GraphRAG dựa trên cơ chế truy vấn động vào đồ thị tri thức kết hợp với sinh ngôn ngữ từ LLM, hành vi của hệ thống không cố định như các API truyền thống. Việc truy xuất tri thức phụ thuộc vào cấu trúc đồ thị, chất lượng embedding, mức độ liên kết giữa các thực thể… khiến quá trình theo dõi lỗi hay xác định điểm nghẽn hiệu năng trở nên khó khăn hơn. Thiếu các công cụ quan sát tiêu chuẩn khiến đội ngũ vận hành khó đánh giá được nguyên nhân sai lệch thông tin lịch sử, xác định tín hiệu bất thường hay tối ưu hiệu suất truy vấn tri thức.

Nếu không có nền tảng triển khai được chuẩn hóa và tích hợp tốt, các hệ thống trả lời câu hỏi lịch sử dựa trên GraphRAG và LLM dễ trở nên phức tạp, khó bảo trì và kém ổn định. Ngược lại, khi các thành phần được thiết kế hợp lý và vận hành bài bản, hệ thống có thể giảm đáng kể gánh nặng kỹ thuật, cải thiện độ tin cậy, nâng cao chất lượng trả lời và mở rộng quy mô linh hoạt. Điều này góp phần thúc đẩy việc ứng dụng GraphRAG vào các hệ thống hỏi – đáp chuyên sâu trong lĩnh vực lịch sử, đáp ứng nhu cầu khai thác tri thức ngày càng tăng trong môi trường giáo dục và nghiên cứu.

\section{Mục tiêu}
Mục tiêu của nghiên cứu này là thiết kế và xây dựng một hệ thống hỏi–đáp lịch sử dựa trên mô hình ngôn ngữ lớn, kết hợp với phương pháp GraphRAG nhằm tăng độ chính xác, khả năng diễn giải và tính tin cậy của câu trả lời. Hệ thống hướng tới khả năng truy xuất tri thức lịch sử có cấu trúc, tự động hóa quá trình xây dựng đồ thị tri thức, tối ưu quá trình truy vấn và đảm bảo vận hành ổn định ở quy mô lớn. \\
Các mục tiêu chính của đề tài bao gồm:

\begin{itemize}
\item \textbf{Xây dựng hệ thống tự động tạo và cập nhật đồ thị tri thức lịch sử:}
Thiết kế cơ chế thu thập, trích xuất thực thể – quan hệ từ tài liệu lịch sử và tự động xây dựng đồ thị tri thức. Hệ thống sử dụng mô hình LLM kết hợp vòng lặp plan–act để phân tích dữ liệu, trích xuất nội dung và liên kết tri thức. Phương pháp Retrieval-Augmented Generation (RAG) được áp dụng để đảm bảo rằng các suy luận và quyết định của mô hình được dựa trên nguồn tư liệu đáng tin cậy.

\item \textbf{Phát triển kiến trúc GraphRAG tối ưu cho bài toán trả lời câu hỏi lịch sử:}
Đề xuất kiến trúc xử lý kết hợp giữa LLM và GraphRAG, cho phép mô hình truy vấn tri thức đồ thị theo cách hiệu quả và nhất quán. Hệ thống được thiết kế theo hướng mô-đun, hỗ trợ mở rộng linh hoạt khi kích thước dữ liệu lịch sử tăng lên. Các kỹ thuật tối ưu như vector indexing, hierarchical retrieval, và local-global graph exploration được tích hợp nhằm tăng tốc độ truy vấn và nâng cao độ chính xác của câu trả lời.

\item \textbf{Xây dựng cơ chế giám sát, đánh giá và phân tích độ tin cậy của câu trả lời:}
Triển khai các công cụ theo dõi hiệu năng truy vấn, độ chính xác thông tin lịch sử, mức độ tương quan giữa câu hỏi và vùng tri thức được truy xuất. Hệ thống hỗ trợ ghi log, phân tích sai lệch (hallucination detection) và truy vết quá trình sinh câu trả lời, giúp đánh giá nguyên nhân sai sót và tăng tính minh bạch cho mô hình.

\item \textbf{Thiết lập quy trình triển khai tự động (CI/CD) cho hệ thống GraphRAG}
Xây dựng pipeline tự động hoá quá trình xử lý dữ liệu, xây dựng đồ thị tri thức, huấn luyện mô hình và triển khai dịch vụ. Tích hợp GitHub Actions hoặc các công cụ tương đương để kiểm thử, cập nhật và phát hành phiên bản mới. Bổ sung cơ chế rollback nhằm đảm bảo hệ thống luôn hoạt động ổn định khi có lỗi trong phiên bản cập nhật.

\end{itemize}

\section{Cấu trúc Khóa luận tốt nghiệp}
Khóa luận được tổ chức thành năm chương chính như sau:

\begin{itemize}
\item \textbf{Chapter 1 – Mở đầu:}
Chương giới thiệu, trình bày bối cảnh và động lực nghiên cứu, mô tả bài toán trả lời câu hỏi lịch sử, nêu rõ mục tiêu và phạm vi của đề tài. Đồng thời, chương cũng giới thiệu tổng quan cấu trúc của khóa luận.

\item \textbf{Chapter 2 – Cơ sở lý thuyết:}
Chương cung cấp tổng quan về các kiến thức nền tảng liên quan đến mô hình ngôn ngữ lớn, phương pháp Retrieval-Augmented Generation, và kỹ thuật GraphRAG. Ngoài ra, chương phân tích các thách thức trong việc xử lý dữ liệu lịch sử, xây dựng đồ thị tri thức và hạn chế của các phương pháp RAG truyền thống.

\item \textbf{Chapter 3 – Phương pháp và Triển khai:}
Chương mô tả chi tiết phương pháp đề xuất, bao gồm quy trình thu thập và tiền xử lý dữ liệu lịch sử, xây dựng đồ thị tri thức, thiết kế kiến trúc GraphRAG, và cơ chế truy vấn tri thức kết hợp LLM. Các mô-đun của hệ thống như trích xuất thực thể – quan hệ, cơ chế lấy ngữ cảnh theo đồ thị, và chiến lược tối ưu truy vấn được trình bày ở đây.

\item \textbf{Chapter 4 – Thực nghiệm và Kết quả:}
Chương trình bày thiết kế thực nghiệm, bộ dữ liệu sử dụng và các tiêu chí đánh giá. Các kết quả thực nghiệm về độ chính xác, mức độ giảm thiểu sai lệch , hiệu suất truy vấn và so sánh giữa LLMs, RAG truyền thống và GraphRAG được báo cáo và phân tích. Chương cũng thảo luận các hạn chế còn tồn tại và đề xuất hướng cải tiến.

\item \textbf{Chapter 5 – Kết luận:}
Chương cuối cùng tóm tắt các đóng góp chính của đề tài, đánh giá tác động của nghiên cứu, và đề xuất các hướng phát triển trong tương lai cho hệ thống hỏi–đáp lịch sử dựa trên LLM và GraphRAG.

\end{itemize}

\newpage